\documentclass[11pt,a4paper]{article}
\usepackage[utf8]{inputenc}
\usepackage{amsmath,amssymb,amsfonts}
\usepackage{geometry}
\usepackage{hyperref}

\geometry{margin=1in}

\title{\textbf{\Huge PolyBLEP Anti-Aliasing Residual Polynomials}\\[0.5em]
\Large Band-Limited Synthesis Math for \texttt{[osc\textasciitilde{}]} and \texttt{[phasor\textasciitilde{}]} Objects}
\author{\textbf{Time Dilation DAW (v0.0.1) Engineering Group}}
\date{\today}

\begin{document}

\maketitle

\begin{abstract}
This document formulates the 2-point PolyBLEP (Polynomial Band-Limited Step) residual function equations utilized in \textbf{Time Dilation DAW} to suppress high-frequency aliasing foldback during non-bandlimited waveform generation.
\end{abstract}

\section{The Aliasing Problem in Naive Digital Oscillators}
A naive digital sawtooth waveform $\phi(t) \in [0, 1)$ generated at normalized frequency $dt = \frac{f}{f_s}$ contains discontinuous step transitions at phase wrap points:
\begin{equation}
x_{\text{naive}}(t) = 2\phi(t) - 1.0
\end{equation}
By Fourier analysis, a ideal continuous sawtooth contains infinite harmonics decreasing as $\frac{1}{k}$. In discrete sample domain, harmonics exceeding the Nyquist limit $f_{\text{Nyquist}} = \frac{f_s}{2}$ fold back into the audible spectrum, causing harsh metallic aliasing noise.

\section{The PolyBLEP Residual Correction Function}
PolyBLEP smooths the sharp step discontinuity by subtracting a 2nd-order polynomial residual $\text{blep}(t, dt)$ within a $\pm 1$ sample neighborhood of the phase wrap transition.

The normalized phase offset $t = \frac{\phi}{dt}$ near transition is evaluated as:
\begin{equation}
\text{blep}(t, dt) =
\begin{cases}
2t - t^2 - 1.0 = -\left(1 - t\right)^2, & 0 \le t < 1 \quad (\text{Just after step}) \\
t^2 + 2t + 1.0 = (t + 1)^2, & -1 < t < 0 \quad (\text{Just before step}) \\
0.0, & \text{otherwise}
\end{cases}
\end{equation}

\section{Application to Waveform Types}

\subsection{Band-Limited Sawtooth Waveform (\texttt{[phasor\textasciitilde{}]}, \texttt{[osc\textasciitilde{}]})}
Subtracting the PolyBLEP residual from the naive sawtooth:
\begin{equation}
y_{\text{saw}}(t) = (2\phi(t) - 1.0) - \text{blep}(\phi(t), dt)
\end{equation}

\subsection{Band-Limited Square / Pulse Waveform (\texttt{[osc\textasciitilde{}]})}
A square wave contains two step discontinuities per period: one at $\phi = 0$ (rising edge) and one at $\phi = 0.5$ (falling edge):
\begin{equation}
y_{\text{square}}(t) = y_{\text{naive}}(t) + \text{blep}(\phi(t), dt) - \text{blep}\left((\phi(t) - 0.5) \pmod{1}, \, dt\right)
\end{equation}

\subsection{Band-Limited Triangle Waveform (\texttt{[osc\textasciitilde{}]})}
Integrating a band-limited square wave yields a smooth triangle wave with attenuation of aliasing by over $80\,\text{dB}$ across all fundamental frequencies up to $f_{\text{Nyquist}}$.

\end{document}
